\documentclass[11pt]{article}
\newcommand{\routelabel}{Intermediate}
\usepackage{eastbridge-handout}
\hypersetup{pdftitle={RPS Workshop: Intermediate}}
\begin{document}
\thispagestyle{plain}

\workshoptitle{Can you beat random play?}{Randomness and prediction}

\begin{context}
Against a bot that always chooses rock, you can win every throw by choosing paper. Against one that chooses each move independently with probability one third, any move you pick wins, loses and draws with equal probability. Its previous moves give you no help with the next one.

\end{context}

We'll calculate payoff as \textbf{+1 for a win, 0 for a draw, -1 for a loss}. Adding those values over a match gives your wins minus your losses. Whoever wins more throws wins the match.


\smallskip\begingroup\small
\setlength{\tabcolsep}{5pt}
\renewcommand{\arraystretch}{1.17}
\begin{tabularx}{\linewidth}{@{}YYYY@{}}
\toprule
\textbf{Your move} & \textbf{Opponent rock} & \textbf{Opponent paper} & \textbf{Opponent scissors} \\
\midrule
Rock & 0 & -1 & +1 \\
Paper & +1 & 0 & -1 \\
Scissors & -1 & +1 & 0 \\
\bottomrule
\end{tabularx}\endgroup\par\smallskip

\subsection*{The Nash equilibrium}

When both players choose independently and uniformly, neither can improve their expected payoff by changing strategy alone. This is a \textbf{Nash equilibrium}. Even a bot with the full match history and a very good prediction algorithm has expected net payoff zero against an independent uniform opponent.


Several house bots do something more predictable. One favors rock; another repeats a sequence; Copycat copies your last move. You can use those habits to predict what comes next and choose a winning reply. Much of the workshop is about learning to recognize them from a short history.


A repeating \code{rock, paper, scissors} bot is an especially helpful example. Count its moves after any complete cycle and you'll find equal totals. But once you've seen rock, you know paper is coming. A random bot can also produce that sequence by chance, so you'll need more than one occurrence before treating it as a pattern.


The kit uses seeded pseudorandom generators to make practice matches repeatable. The exercises use those seeds to compare versions under the same conditions. They assume each bot's current random choice is private; reconstructing an opponent's generator from its seed would be a different problem.
\par

\clearpage
\workshoptitle{What your bot sees}{Inside bot.py}

\begin{context}
For each match, the arena starts a fresh Python process and imports bot.py. It calls \code{setup(config)} once, if you've defined it, then asks \code{next\_move(...)} for each throw. Here is the starter bot:

\end{context}

\begin{lstlisting}[style=python]
from random import Random
from rpsdk import Move

_rng = Random()

def setup(config):
    _rng.seed(config["seed"])

def next_move(my_history, opponent_history, match_state):
    return _rng.choice(list(Move))
\end{lstlisting}

Return \code{Move.ROCK}, \code{Move.PAPER} or \code{Move.SCISSORS} from next\_move. These values also have a comparison method: \code{a.beats(b)} tells you whether a beats b. The arena accepts lowercase strings too, though the named values help catch spelling mistakes.


\subsection*{Reading the histories}

Both history lists contain \textbf{completed throws, oldest first}. On the first call, both are empty. After two throws, you might see:


\begin{lstlisting}[style=python,numbers=none,xleftmargin=4pt,framexleftmargin=0pt]
my_history       == [Move.ROCK, Move.PAPER]
opponent_history == [Move.SCISSORS, Move.ROCK]
\end{lstlisting}

You won both of those throws. \code{opponent\_history[-1]} gives their last move, while \code{my\_history[-1]} gives yours. Check that a list has an entry before using \code{[-1]}, or two entries before using \code{[-2]}. The histories never include the throw you're about to choose.


\smallskip\begingroup\small
\setlength{\tabcolsep}{5pt}
\renewcommand{\arraystretch}{1.17}
\begin{tabularx}{\linewidth}{@{}L{.28\linewidth}Y@{}}
\toprule
\textbf{State entry} & \textbf{Meaning} \\
\midrule
\code{round} & Zero-based throw index. The next call above has round 2. \\
\code{best\_of} & Total throws in this match, including draws. \\
\code{last\_outcome} & Your preceding throw's win, loss or draw; None initially. \\
\code{seed} & Your bot's reproducibility seed. \\
\code{timeouts}, \code{opponent\_timeouts} & Accumulated server timeouts in this match. \\
\bottomrule
\end{tabularx}\endgroup\par\smallskip

Each match starts with empty histories and fresh learned state. Reset any global counters in setup so your bot also works when tests call setup repeatedly. Seed the random generator there once; next\_move can then draw from it throughout the match.


\begin{checkpoint}
\textbf{Try it on paper:} Copycat copies your preceding move. Given the two histories above, what will it play next? Which history did you use?
\end{checkpoint}

\clearpage
\workshoptitle{Count the moves}{Route B / Intermediate}

\begin{context}
Biased Random plays rock 60\% of the time, paper 25\% and scissors 15\%. Against it, paper has expected net payoff $0.60-0.15=+0.45$ per throw. We'll build a bot that estimates those probabilities from the history, so it can respond to biases we haven't told it about.

\end{context}

\lessonstep{Step 1}{count what you have seen}

Use \code{collections.Counter} to count moves in opponent\_history. Choose a random move while the history is empty. Then build a response from those counts:


\begin{lstlisting}[style=python]
from collections import Counter

counts = Counter(opponent_history)
values = {
    Move.ROCK: counts[Move.SCISSORS] - counts[Move.PAPER],
    Move.PAPER: counts[Move.ROCK] - counts[Move.SCISSORS],
    Move.SCISSORS: counts[Move.PAPER] - counts[Move.ROCK],
}
\end{lstlisting}

Each value counts how many of the observed moves it would have beaten, minus how many would have beaten it. Divide by the history length to get an estimated payoff per throw. Either way, the largest value identifies the same move. Choose randomly among ties.


Do not automatically counter the most frequent move. If the distribution is \textbf{40\% rock, 21\% paper, 39\% scissors}, paper's expected payoff is only $+0.01$, while rock's is $+0.18$. Rock exploits the abundant scissors without losing often to paper.


\begin{checkpoint}
\textbf{Try the calculation:} work out all three expected payoffs for the 40/21/39 distribution using the payoff table. Why does rock have a higher expected payoff than paper?
\end{checkpoint}

\lessonstep{Step 2}{test the estimates}

\begin{lstlisting}[style=shell,numbers=none,xleftmargin=4pt,framexleftmargin=0pt]
rps-cli play --against biased_random --games 5 --seed 100
rps-cli play --against random_uniform --games 5 --seed 100
\end{lstlisting}

You should have an advantage against Biased Random. Against Random Uniform, your counts will still be uneven, particularly early in a match, but they won't predict its next move. Try several seeds and watch how the results vary.


Try waiting for 5, 15 or 30 observations before using the counts, playing randomly until then. Keep the rest of your code unchanged. Waiting gives you more data to estimate from, at the cost of some throws you could have won against Biased Random.
\par
\clearpage
\section*{An opponent that changes its mind}

\begin{context}
Switcheroo plays rock for the first half of the series and scissors afterwards. By the time it switches, your bot has counted hundreds of rocks. Those old observations can keep it choosing paper for the rest of the match.

\end{context}

\lessonstep{Step 3}{shorten the memory}

Start with a rolling window:


\begin{lstlisting}[style=python,numbers=none,xleftmargin=4pt,framexleftmargin=0pt]
window = opponent_history[-30:]
counts = Counter(window)
\end{lstlisting}

The payoff calculation is unchanged. Compare windows of 10, 30 and 100. A short window reacts quickly but has more sampling noise. A long window estimates a fixed bias more smoothly but carries old behavior across a change.


\begin{lstlisting}[style=shell,numbers=none,xleftmargin=4pt,framexleftmargin=0pt]
rps-cli play --against switcheroo,shifting_bias --games 5 --seed 200
rps-cli play --against biased_random --games 5 --seed 200
\end{lstlisting}

Shifting Bias changes its favored move every 50 throws; the favorite has probability 70\%, and the others 15\% each. Its initial favorite varies with the seed. For each window size, count how many throws your bot needs to recover after a switch.


\lessonstep{Step 4}{inspect the change}

\begin{lstlisting}[style=shell,numbers=none,xleftmargin=4pt,framexleftmargin=0pt]
rps-cli play --against shifting_bias --seed 200 --output drift.json
\end{lstlisting}

The JSON file contains a list of matches, each with its seed, totals and every observed throw. Each throw has a zero-based \code{round}, your move, the opposing move, a +1/0/-1 \code{outcome}, and any local error.


Read the first match in a small Python script:


\begin{lstlisting}[style=python]
import json
from pathlib import Path
match = json.loads(Path("drift.json").read_text())["matches"][0]
for start in range(0, len(match["rounds"]), 50):
    block = match["rounds"][start:start + 50]
    print(start, sum(row["outcome"] for row in block))
\end{lstlisting}

In a 501-throw match, the final block has just one throw. Divide each block's total by its length if you want to compare average payoffs.


\begin{checkpoint}
\textbf{Compare:} which window worked best against Switcheroo? Did the same window work best against Biased Random? Look at the throws around a switch to explain any difference.
\end{checkpoint}
\clearpage
\section*{Using the previous throw}

\begin{context}
Sticky is balanced in the long run, yet it repeats its previous move with probability 0.8. The other two moves each have probability 0.1. If its last move was rock, paper has expected net payoff $0.8-0.1=+0.7$ on the next throw.

\end{context}

\lessonstep{Step 5}{respond to Sticky}

Try countering Sticky's last move and compare the result with your counting bot. Overall counts mix together the throws after rock, paper and scissors. Keeping those cases separate lets you use Sticky's tendency to repeat.


Now inspect \textbf{Win Stay Lose Shift}. It repeats after a win \textbf{or a draw}, and after a loss chooses uniformly between the two other moves. You can reconstruct its previous outcome from the histories:


\begin{lstlisting}[style=python,numbers=none,xleftmargin=4pt,framexleftmargin=0pt]
opponent_won = opponent_history[-1].beats(my_history[-1])
drawn = opponent_history[-1] == my_history[-1]
\end{lstlisting}

These lines require a nonempty history. \code{match\_state["last\_outcome"]} is from \textbf{your} perspective, so your win means its loss.


If the opponent just lost with rock, it will next choose paper or scissors. Your best reply is scissors: it wins against paper and draws against scissors. Playing rock would win once and lose once on average. Work out the corresponding replies after losses with paper and scissors.


\lessonstep{Step 6}{predict a reactive opponent}

After its opening, Contrarian always counters your previous move. If you played rock, it will choose paper, so your reply should be scissors. Use \textbf{my\_history} to make this prediction.


\begin{lstlisting}[style=shell,numbers=none,xleftmargin=4pt,framexleftmargin=0pt]
rps-cli play --against sticky,win_stay_lose_shift,contrarian
\end{lstlisting}

Your tournament bot will have to decide which rule fits the opponent. Count how many recent observations agree with each rule, and use the random fallback when none fits well.


\begin{checkpoint}
\textbf{Check:} write the prediction rule for Sticky, Win Stay Lose Shift and Contrarian. Walk through an opening, a win and a loss for each. Does your code use the right player's history and outcome?
\end{checkpoint}

\textbf{Next route:} learn Sticky's probabilities from three separate tables, one for each last opposing move. This is a first-order Markov model.
\par
\clearpage
\section*{Comparing your bots}

\begin{context}
You now have several strategies to compare: counts over the whole history, a short window, and rules that respond to the previous throw. Save each as a separate version and run them against the same opponents.

\end{context}

\lessonstep{Step 7}{reserve a test set}

Use one group of seeds while designing your strategy and another group for the final comparison. For example, tune on 100 through 104, then test on 900 through 904. Keep the opponents and series length the same for both versions.


\begin{lstlisting}[style=shell,numbers=none,xleftmargin=4pt,framexleftmargin=0pt]
rps-cli play --against biased_random,shifting_bias,sticky \
  --best-of 501 --games 5 --seed 900 --output test-set.json
\end{lstlisting}

Every match starts with a fresh process and fresh learned state. You cannot carry a table learned against one opponent into the next match. If your bot uses global counters, reset them in setup so unit tests and other local uses also start correctly.


\smallskip\begingroup\small
\setlength{\tabcolsep}{5pt}
\renewcommand{\arraystretch}{1.17}
\begin{tabularx}{\linewidth}{@{}L{.25\linewidth}YYYY@{}}
\toprule
\textbf{Policy} & \textbf{Fixed bias} & \textbf{Changing bias} & \textbf{Reactive bot} & \textbf{Uniform control} \\
\midrule
\rule{0pt}{22pt}All-history counts &  &  &  &  \\
\rule{0pt}{22pt}Recent window &  &  &  &  \\
\rule{0pt}{22pt}Your combined policy &  &  &  &  \\
\bottomrule
\end{tabularx}\endgroup\par\smallskip

For each column, record \textbf{match W-L-D} and average \textbf{net payoff per throw}. Keep the results for each opponent separate. Otherwise, a large margin against Rocky can hide a string of narrow losses against the rest of the field.


\lessonstep{Step 8}{mix in random moves}

Try using your learned policy on 80\% of throws and choosing uniformly at random on the rest. This gives an opponent fewer predictable moves to respond to. You'll also win fewer throws where your prediction was already correct. Add the mixed version to your comparison.


Each time you choose a window, threshold or mixture based on its score, you're tuning to those matches. Use a new batch of seeds for the final comparison to see how well the choice carries over.


\subsection*{Submit your latest version}

Run \code{rps-cli test} and \code{rps-cli validate}, then submit under your existing team name and check that the latest version is active. Keep your earlier working bot available in case an ambitious change introduces an error near the end.


\begin{checkpoint}
\textbf{Discuss:} show a partner a match where your bot struggled. What was it predicting? How many throws did it take to change its prediction, and would a different window have helped?
\end{checkpoint}
\end{document}
